\section{Homomorphisms of Free Modules, I}
\subsection{Matrices}
As pointed out previously, we can classify free modules by dimension; that is, for all free modules $F$, there is a set $A$ unique up to isomorphism such that $F\cong R^{\oplus A} $. This means we can understand $\Hom_{R}(F_1, F_2) $ by simply understanding $\Hom_{R}(R^{\oplus A_1} , R^{\oplus A_2} ) $, but unfortunately this identification is not canonical, as it depends on our choice of isomorphism $F_1\cong R^{\oplus A_1} $. That is, our representation depends on our choice of basis. We will deal with this ambiguity later; for now, we will attempt to describe $\Hom_{R}(R^{\oplus A_1} , R^{\oplus A_2} ) $. This is fairly simple when our modules are finitely generated, so we will focus on this case for now.


